\section{Conclusion}
\label{sec:conclusion}

This paper addressed a mismatch between what private hierarchical federated
learning guarantees and what clinical IoT requires: existing systems protect
individual records, whereas a clinical deployment must protect patients, each
of whom contributes many records. The proposed framework, PP-HFM, closes this
gap by construction. Each client sums a patient's gradients before clipping,
calibrates its noise to its own population, which parallel composition
permits for every patient held by one client, and accounts for every
release of the hierarchy, edge tier included. The guarantee was proved for the
training transcript; it holds end to end only when the encoders and
standardization constants are fitted on public data, which on the evaluated
corpora is the eICU configuration, and there for the $97.1\,\%$ of training
individuals treated at a single hospital.

The evaluation on three corpora supports four conclusions. First, the system
objectives are met on every corpus with two exceptions on eICU and one on
WEAR, where the record-level utility loss of $1.97$~pp is inconclusive
against the two-point target: the training
budget is at most $\varepsilon = 0.9$ ($1.0$ with the evaluated conformal
release), uplink traffic falls by $90.8$--$96.1\,\%$, modeled energy per
round by $72.3$--$79.6\,\%$, and the $95$th-percentile inference latency on
a workstation CPU is at most $2.3$~ms, but on eICU the utility loss exceeds
the two-point target end to end and is inconclusive against it with
training-fitted constants, and update delivery reaches
$99\,\%$ only from eight edges. Second, the patient-level guarantee is affordable: it costs no utility on
VitalDB, partitioned by person, and on eICU it costs $3.1$ AUROC points $[1.4, 4.8]$ ($3.9$ $[2.4, 5.4]$ end to
end), a cost set by the smallest hospitals. Third, on VitalDB and WEAR, robustness to a strict
budget is determined by the initialization, not by the size of the privatized
subset: pretrained initialization is worth $0.199$ AUROC at
$\varepsilon = 0.5$, whereas restricting noise to the head changes accuracy by
at most $0.006$ under privacy and is justified by communication alone. Fourth, the price of
an individual-level guarantee follows $A(1 - 1/k)$ in records per individual
and vanishes where the two units coincide, a scaling that holds out of sample
on the first corpus and in shape on the second. On VitalDB, a patient-level
audit with a positive control is consistent with the distinction: record-level
training remains detectably above its difficulty floor under offline LiRA,
whereas patient-level training does not; the audit bounds these attacks
without establishing the absence of leakage. The corpora carry different
evidence: eICU tests the guarantee end to end for individuals at one hospital, VitalDB tests it on the
training transcript, and WEAR, where each client is one participant,
evaluates the system at record level.

Several directions remain open. First, a third corpus is needed to determine
whether the ceiling of the cost curve is a constant of the guarantee or a
fraction of the precision attainable on a task. Second, the noise imposed by
the smallest clients could be reduced, for example by sampling patients across
the clients of an edge. Third, clipping around an estimate of each
individual's mean gradient, which the analysis indicates could recover the
cost of the unit, should be tested. Finally, the audit should be extended to eICU
and, with a larger shadow population, to low false-positive rates.
